%% The Motivic Shadow Tower.
%%
%% Constructs the motivic lift of the Vol I shadow tower on affine
%% Arnold screens. The motivic object is weight-filtered, not purely
%% weight r in general: rational Virasoro coefficients live in the
%% Tate-scalar lane, while non-rational Arnold periods live in the
%% Brown--Goncharov motivic MZV filtration.
%%
%% Vol I theory sequence. Depends on shadow tower (see \ref{part:characteristic-datum}), Vol II
%% GRT-parametrized seven faces, Vol III shadow-Feynman
%% identification S_r = delta^{(r)}.
%%
%% All macros \providecommand-guarded.

\providecommand{\MZV}{\mathrm{MZV}}
\providecommand{\mot}{\mathrm{mot}}
\providecommand{\num}{\mathrm{num}}
\providecommand{\GRT}{\mathrm{GRT}_1}
\providecommand{\grt}{\mathfrak{grt}_1}
\providecommand{\Face}{\mathrm{Face}}
\providecommand{\cA}{\mathcal{A}}
\providecommand{\cF}{\mathcal{F}}
\providecommand{\PhiFA}{\Phi^{\mathrm{FA}}}
\providecommand{\SpCh}{\mathrm{Sp}^{\mathrm{ch}}}
\providecommand{\cW}{\mathcal{W}}
\providecommand{\Walg}{\mathcal{W}}
\providecommand{\Vir}{\mathrm{Vir}}
\providecommand{\MC}{\mathrm{MC}}
\providecommand{\wt}{\mathrm{wt}}
\providecommand{\rat}{\mathrm{rat}}
\providecommand{\KZ}{\mathrm{KZ}}
\providecommand{\Delta}{\Delta}
\providecommand{\zet}{\zeta}
\providecommand{\Li}{\mathrm{Li}}
\providecommand{\per}{\mathrm{per}}
\providecommand{\coll}{\mathrm{coll}}
\providecommand{\Res}{\operatorname{Res}}
\providecommand{\fg}{\mathfrak{g}}
\providecommand{\Fil}{\mathrm{Fil}}
\providecommand{\ClaimStatusConditional}{}

\chapter{The Motivic Shadow Tower}
\label{ch:motivic-shadow-tower}

%%%%%%%%%%%%%%%%%%%%%%%%%%%%%%%%%%%%%%%%%%%%%%%%%%%%%%%%%%%%%%%%%%%%%%%%%%%%%
\section*{Prologue}
%%%%%%%%%%%%%%%%%%%%%%%%%%%%%%%%%%%%%%%%%%%%%%%%%%%%%%%%%%%%%%%%%%%%%%%%%%%%%

The shadow tower $\{S_r(\cA)\}_{r \geq 2}$ is the sequence of scalar
coefficients obtained from the ordered Maurer--Cartan element
$\Theta_{\cA}$ after Arnold-residue extraction and symmetric averaging.
The low Virasoro values
\[
 S_2(\Vir_c)=\frac c2,\qquad S_3(\Vir_c)=2,\qquad
 S_4(\Vir_c)=\frac{10}{c(5c+22)},\qquad
 S_5(\Vir_c)=\frac{-48}{c^2(5c+22)}
\]
are rational functions of the central charge. Their rationality is not
a motivic purity statement in weight $r$: it says that these particular
Arnold pairings have only the Tate-scalar component. Higher-spin
channels may carry genuine Brown--Goncharov periods, but a non-zero
period coefficient must be produced by an explicit Arnold/KZ iterated
integral that survives the ordered-to-symmetric averaging map.

The GRT-parametrised faces act on the same Arnold/KZ period algebra.
Thus the natural motivic object attached to a shadow coefficient is a
weight-filtered motivic period, together with its rational OPE
coefficient, not a bare numerical scalar.

Throughout, we distinguish numerical multiple zeta values
$\MZV^{\num} \subset \mathbb{R}$ (the $\mathbb{Q}$-algebra generated
by $\zet(n_1, \ldots, n_k)$ with $n_1 \geq 2$) from motivic
multiple zeta values $\MZV^{\mot}$ (Goncharov--Brown
$\mathbb{Q}$-algebra of motivic periods of the mixed Tate motive
$\mathrm{MT}(\mathbb{Z})$). The period map
$\per \colon \MZV^{\mot} \to \MZV^{\num}$ is surjective (Brown
2012); injectivity is the motivic period conjecture in this mixed Tate
setting, not a theorem. The motivic ring $\MZV^{\mot}$ carries a
coaction of the motivic Galois group of
$\mathrm{MT}(\mathbb{Z})$. Through the Ihara--Drinfeld action on the
motivic fundamental group of
$\mathbb P^1\setminus\{0,1,\infty\}$ this coaction is compatible with
the Grothendieck--Teichmuller Lie algebra. No argument below uses an
identification of the full motivic Galois Lie algebra with $\grt$.

Three separations govern the chapter. A shadow coefficient has an
external bar weight $r$, while its motivic period may have any Brown
weight $\leq r$; rational Virasoro coefficients have Brown weight $0$.
Arnold forms control affine ordered screens; projective and
positive-genus curves add the extra relations recorded in
Appendix~\ref{app:arnold-relations}. Finally, shadow/Gel'fand--Fuks
representatives are not curve-level chiral Hochschild classes.

%%%%%%%%%%%%%%%%%%%%%%%%%%%%%%%%%%%%%%%%%%%%%%%%%%%%%%%%%%%%%%%%%%%%%%%%%%%%%
\section{Setup: motivic multiple zeta values and the weight filtration}
\label{sec:mot-setup}
%%%%%%%%%%%%%%%%%%%%%%%%%%%%%%%%%%%%%%%%%%%%%%%%%%%%%%%%%%%%%%%%%%%%%%%%%%%%%

We recall only what is required; the reader is referred to
Brown's \emph{Mixed Tate motives over $\mathbb{Z}$}
(arXiv:1102.1312) and Deligne--Goncharov's \emph{Groupes
fondamentaux motiviques de Tate mixte} (Ann. Sci. \'Ecole Norm.
Sup. 38, 2005) for the foundational statements. The proofs in this
chapter use the ring $\MZV^{\mot}$ and its coaction through the cited
theorems stated inline.

\begin{definition}[\label{def:mzv-mot}Motivic multiple zeta values]
Let $\mathrm{MT}(\mathbb{Z})$ denote the Tannakian category of
mixed Tate motives over $\mathbb{Z}$ (Deligne--Goncharov 2005,
Brown 2012). Write $\MZV^{\mot}$ for the $\mathbb{Q}$-algebra of
\emph{motivic periods} of $\mathrm{MT}(\mathbb{Z})$: for each
collection of indices $(n_1, \ldots, n_k)$ with $n_1 \geq 2$ and
$n_i \geq 1$, there is a distinguished element
$\zet^{\mot}(n_1, \ldots, n_k) \in \MZV^{\mot}$; the weight is
$\wt = n_1 + \cdots + n_k$. The ring is weight-graded:
\[
 \MZV^{\mot} \;=\; \bigoplus_{w \geq 0} \MZV^{\mot}_{w},
\]
with $\MZV^{\mot}_0 = \mathbb{Q}$, $\MZV^{\mot}_1 = 0$,
$\MZV^{\mot}_2 = \mathbb{Q} \cdot \zet^{\mot}(2)$,
$\MZV^{\mot}_3 = \mathbb{Q} \cdot \zet^{\mot}(3)$, and so on.
\end{definition}

\begin{remark}[\label{rem:wt-1-vanishes}No motivic multiple zeta
value of weight one]
Definition~\ref{def:mzv-mot} entails $\MZV^{\mot}_1 = 0$: any
multiple zeta value of weight $1$ would be proportional to
$\zet(1)$, which diverges. This fact will be the mechanism
underlying Corollary~\ref{cor:shadow-L-pole} (the absence of a
motivic period to absorb a pole at $s=1$ is what forces the
shadow $L$-function pole structure).
\end{remark}

\begin{definition}[\label{def:period-map}Period map]
The \emph{period map}
\[
 \per \colon \MZV^{\mot} \longrightarrow \mathbb{R}
\]
is the $\mathbb{Q}$-algebra homomorphism sending
$\zet^{\mot}(n_1, \ldots, n_k)$ to the numerical multiple zeta
value $\zet(n_1, \ldots, n_k) \in \mathbb{R}$. The image is
$\MZV^{\num} \subset \mathbb{R}$. By Brown's period theorem
(\emph{loc.\ cit.}\ Thm.\ 1.1), $\per$ is surjective onto
$\MZV^{\num}$.
\end{definition}

\begin{definition}[\label{def:mot-coaction}Motivic coaction]
The ring $\MZV^{\mot}$ is a comodule over the coordinate Hopf algebra
\[
 \mathcal{G}^{\mot}
 =
 \mathcal{O}\bigl(\mathrm{Gal}^{\mot}(\mathrm{MT}(\mathbb{Z}))\bigr).
\]
Brown identifies the Lie algebra of the unipotent radical, after a
choice of de~Rham generators, with the free graded Lie algebra having
one generator $\sigma_{2n+1}$ in each odd weight $2n+1\ge3$ (Brown
2012 Thm.~1.2). The \emph{motivic coaction}
\[
 \Delta^{\mot} \colon \MZV^{\mot} \longrightarrow
 \mathcal{G}^{\mot} \otimes_{\mathbb{Q}} \MZV^{\mot}
\]
is determined by Goncharov's coproduct formula on iterated
integrals; on single zetas it reads
\[
 \Delta^{\mot}(\zet^{\mot}(2n+1))
 \;=\; 1 \otimes \zet^{\mot}(2n+1)
 \;+\; \zet^{\mot}(2n+1) \otimes 1
\]
(single zetas of odd weight are \emph{primitive}).
\end{definition}

The Ihara--Drinfeld action sends the motivic Galois Lie algebra to
the Grothendieck--Teichm\"uller Lie algebra acting on KZ associators.
We use the notation $\sigma_{2n+1}$ for this motivic image in $\grt$.
Nothing in the sequel requires the unproved assertion that the image is
all of $\grt$. The depth and dimension conjectures of
Zagier--Hoffman and Broadhurst--Kreimer are likewise not inputs; the
arguments use only Brown's motivic coaction, the weight filtration, and
the explicit low-weight lines $\MZV^{\mot}_0$, $\MZV^{\mot}_1$, and
$\MZV^{\mot}_3$.

\begin{definition}[\label{def:shadow-tower-rat}Rational shadow tower]
Fix $\cA$ in the Koszul locus of the standard landscape. Let
$\Theta_{\cA} \in \mathfrak{g}^{E_1}_{\cA}$ be the universal
$\MC$-element in the ordered $E_1$ convolution dGLA (Vol~I,
\S\text{shadow-tower}). Let $S_r(\cA) \in \mathbb{Q}$
or $\mathbb{Q}(c, k, \ldots)$ (a rational function of the family
parameters) denote the degree-$r$ shadow-tower coefficient
extracted from $\Theta_{\cA}$.

The \emph{period-scalar shadow ring} is
\[
 R^{\rat}(\cA) \;=\; \mathbb{Q}\bigl[(2\pi i)^{2n}, \zet(2m+1)
 \mid n\geq1,\ m\geq1\bigr]
\]
(the $\mathbb{Q}$-subalgebra of $\mathbb{C}$ generated by the even
Tate periods and the odd single zetas), graded by period weight with
$\deg((2\pi i)^{2n}) = 2n$ and
$\deg(\zet(2m+1)) = 2m+1$.
\end{definition}

\begin{remark}[\label{rem:even-zeta-2pi-i}Even zetas reduce to
powers of $2 \pi i$]
$\zet(2n) = -\tfrac{(2\pi i)^{2n}}{2 (2n)!} B_{2n}$ where
$B_{2n}$ is a Bernoulli number; thus all even zetas already lie
in $\mathbb{Q}[(2\pi i)^{2n}]$. The choice to generate $R^{\rat}$
by $(2\pi i)^2$ and odd zetas is economical: it matches the
generators of $\MZV^{\mot}$, with $\zet^{\mot}(2)$ and
$(2\pi i)^{\mot,2}$ identified via
$\zet^{\mot}(2) = -\tfrac{(2\pi i)^{\mot,2}}{24}$.
\end{remark}

%%%%%%%%%%%%%%%%%%%%%%%%%%%%%%%%%%%%%%%%%%%%%%%%%%%%%%%%%%%%%%%%%%%%%%%%%%%%%
\section{The motivic-lift theorem}
\label{sec:mot-lift-thm}
%%%%%%%%%%%%%%%%%%%%%%%%%%%%%%%%%%%%%%%%%%%%%%%%%%%%%%%%%%%%%%%%%%%%%%%%%%%%%

\begin{theorem}[\label{thm:shadow-tower-motivic-lift}Motivic lift of
Arnold-period shadow coefficients]
\ClaimStatusProvedHere
Let $\cA$ lie in the Koszul locus of the standard landscape. Fix the
genus-zero affine Arnold/KZ model with tangential base points. Suppose
the ordered residue computation of the $r$-th shadow coefficient is
written as a finite expression
\[
 S_r(\cA)=\sum_{\nu} q_{\nu,r}(\cA)\,\per(I_{\nu,r}^{\mot}),
 \qquad q_{\nu,r}(\cA)\in\mathbb Q(c,k,\ldots),
\]
where each $I_{\nu,r}^{\mot}$ is a motivic iterated integral of Arnold
length at most $r$. Then
\[
 S_r^{\mot}(\cA):=\sum_{\nu} q_{\nu,r}(\cA)\,I_{\nu,r}^{\mot}
 \;\in\;
 \Fil^W_{\le r}\MZV^{\mot}\otimes_{\mathbb Q}\mathbb Q(c,k,\ldots)
\]
is a motivic lift and satisfies
\[
 \per\bigl(S_r^{\mot}(\cA)\bigr)=S_r(\cA).
\]
If the Arnold expression has only the unit period, then
$S_r^{\mot}(\cA)\in \MZV^{\mot}_0\otimes\mathbb Q(c,k,\ldots)$; its
shadow degree remains $r$, but its Brown motivic weight is $0$.
The construction is functorial for morphisms preserving the chosen
ordered residue model and is compatible with rational
ordered-to-symmetric averaging.
\end{theorem}

\begin{proof}
Arnold's theorem identifies the cohomology of the affine configuration
space with the Orlik--Solomon algebra generated by the logarithmic
forms $\eta_{ij}=d\log(z_i-z_j)$ modulo the triangular circuit
relations. Deligne--Goncharov and Brown lift the corresponding
iterated integrals on the genus-zero KZ/Arnold model to
$\MZV^{\mot}$, with period map equal to the numerical iterated
integral.

The ordered Maurer--Cartan element has the form
\[
 \Theta_{\cA} \;=\; r(z) \,dz
 \;+\; \sum_{r \geq 3} \theta_r \cdot \omega_r,
\]
where the $\omega_r$ are Arnold words on affine tangent screens and
the OPE coefficients in $\theta_r$ are rational functions of the
family parameters. The scalar shadow coefficient is obtained by
applying the rational averaging map
\[
 \mathrm{av}^{\mathrm{mod}}
 \colon \mathfrak{g}^{E_1} \longrightarrow \mathfrak{g}^{\mathrm{mod}}.
\]
Since averaging is $\mathbb Q$-linear and acts on the ordered
configuration labels rather than on the period algebra, it commutes
with replacing each numerical Arnold iterated integral by its motivic
lift. Arnold length bounds Brown weight, so the resulting coefficient
lies in $\Fil^W_{\le r}\MZV^{\mot}$. Applying $\per$ termwise gives
the original scalar. Functoriality is the functoriality of the OPE
coefficient system and of the Arnold/KZ iterated-integral torsor under
maps that preserve the chosen ordered residue model.
\end{proof}

\begin{remark}[\label{rem:mot-lift-independence}Associator torsor]
Changing the motivic associator changes the coordinates on the same
motivic torsor. The intrinsic assertion is therefore comodule
stability, not equality of coordinate formulas. A displayed expression
for $S_r^{\mot}$ is tied to the chosen Arnold/KZ coordinates; its
period and its motivic coaction orbit are coordinate-independent.
\end{remark}

%%%%%%%%%%%%%%%%%%%%%%%%%%%%%%%%%%%%%%%%%%%%%%%%%%%%%%%%%%%%%%%%%%%%%%%%%%%%%
\section{The GRT motivic coaction on the shadow tower}
\label{sec:grt-coaction}
%%%%%%%%%%%%%%%%%%%%%%%%%%%%%%%%%%%%%%%%%%%%%%%%%%%%%%%%%%%%%%%%%%%%%%%%%%%%%

Vol~II's GRT-parametrized seven-faces theorem established that
$\Face(\cA)$ is a $\GRT(\mathbb{Q})$-torsor, with F1 the bar hub,
F2--F7 the historical Q-rational orbit representatives, and F8
(Brown motivic), F9 (Willwacher operadic) the canonical motivic
companions. The action of $\GRT(\mathbb{Q})$ on r-matrices is
the associator transport; on Arnold period coordinates it is induced
by the motivic coaction below.

\begin{theorem}[\label{thm:grt-motivic-coaction}Motivic coaction on
the Arnold shadow-period envelope]
\ClaimStatusProvedHere
For fixed $\cA$, $r$, and ordered residue expressions as in
Theorem~\ref{thm:shadow-tower-motivic-lift}, let
\[
 \mathscr S^{\mot}_{\le r}(\cA)
 \subset
 \Fil^W_{\le r}\MZV^{\mot}\otimes_{\mathbb Q}\mathbb Q(c,k,\ldots)
\]
be the $\mathbb Q(c,k,\ldots)$-span of all motivic Arnold iterated
integrals of length at most $r$ that occur either in the expressions
$S_j^{\mot}(\cA)$, $2\le j\le r$, or as right tensor factors in an
iterated Goncharov coproduct of such integrals. For $b\le r$, let
$\mathscr S^{\mot}_{\le b}(\cA)$ be the length-$\le b$ part of this
same envelope. Then $\mathscr S^{\mot}_{\le r}(\cA)$ is a
$\mathcal G^{\mot}$-subcomodule:
\[
 \Delta^{\mot}\bigl(\mathscr S^{\mot}_{\le r}(\cA)\bigr)
 \subset
 \bigoplus_{a+b\le r}
 \mathcal G^{\mot}_{a}\otimes
 \mathscr S^{\mot}_{\le b}(\cA).
\]
In any homogeneous Brown basis one may write
\[
 S_r^{\mot}(\cA)=\sum_{b\le r} S_{r,b}^{\mot}(\cA),
 \qquad S_{r,b}^{\mot}(\cA)\in \MZV^{\mot}_{b}\otimes\mathbb Q(c,k,\ldots),
\]
and the coaction lowers the right tensor weight by the left tensor
weight. The formula is basis-dependent; the subcomodule statement is
intrinsic. The smaller literal span of the periods appearing in the
displayed coefficients is not asserted to be stable unless it already
contains these coproduct right factors.
\end{theorem}

\begin{proof}
Goncharov's coproduct for motivic iterated integrals is deconcatenation
with admissible subintegrals. A subword of an Arnold word on an affine
collision screen is again an Arnold word of smaller or equal length,
and the Orlik--Solomon relations are rational. Therefore Brown's
coaction sends every motivic period appearing in a shadow coefficient
to a finite sum of left motivic Galois coefficients tensor Arnold
periods of smaller Brown weight, and those right tensor factors lie in
the envelope by definition.

The averaged motivic map
\[
 \mathrm{av}^{\mot}:=\mathrm{av}^{\mathrm{mod}}\otimes_{\mathbb Q}
 \mathrm{id}_{\MZV^{\mot}}
\]
is rational on the configuration labels, so it commutes with the
coaction on the coefficient ring. The weight bound follows because a
length-$r$ Arnold word has no motivic subperiod of Brown weight
exceeding $r$. The GRT torsor acts by changing associator coordinates;
the subcomodule does not depend on that coordinate choice.
\end{proof}

\begin{remark}[\label{rem:f8-f9-coincide}$F_8$ and $F_9$ give the
same coaction]
The Willwacher operadic and Brown motivic realisations act through the
same motivic image in the $\grt$-torsor after the Vol~II comparison.
The intrinsic object is the coaction envelope of
Theorem~\ref{thm:grt-motivic-coaction}; coordinates coming from $F_8$
or $F_9$ give different names to the same orbit.
\end{remark}

%%%%%%%%%%%%%%%%%%%%%%%%%%%%%%%%%%%%%%%%%%%%%%%%%%%%%%%%%%%%%%%%%%%%%%%%%%%%%
\section{Explicit motivic computations}
\label{sec:motivic-shadow-explicit-computations}
%%%%%%%%%%%%%%%%%%%%%%%%%%%%%%%%%%%%%%%%%%%%%%%%%%%%%%%%%%%%%%%%%%%%%%%%%%%%%

We compute two explicit Virasoro motivic shadow coefficients and then
isolate the first higher-spin period boundary. The Virasoro lifts are
Tate-scalar: no $\zet^{\mot}$ enters. The $W_3$ weight-$6$ statement
does not prove a non-zero $\zet^{\mot}$ coefficient; it names the
Arnold/KZ pairing that must be computed before such a coefficient can
be asserted.

\subsection{$S_4^{\mot}(\Vir_c)$}

\begin{proposition}[\label{prop:s4-vir-mot}Motivic lift of
$S_4(\Vir_c)$]
\ClaimStatusProvedHere
The motivic lift of the weight-$4$ Virasoro shadow coefficient is
\[
 S_4^{\mot}(\Vir_c) \;=\; \frac{10}{c (5c + 22)}
 \;\in\; \MZV^{\mot}_{0} \otimes \mathbb{Q}(c)
 \;\subset\; \Fil^W_{\le 4}\MZV^{\mot} \otimes \mathbb{Q}(c).
\]
The external shadow degree is $4$; the Brown motivic weight is $0$.
\end{proposition}

\begin{proof}
The landscape census fixes the Virasoro coefficient on the nonsingular
surface $c(5c+22)\ne0$:
\[
 S_4(\Vir_c) \;=\; \frac{10}{c (5c + 22)}.
\]
The denominator factors as
$\det G_{\wt\,=\,2}(\Vir_c) = c^2(5c+22)/2$, the weight-$2$
Shapovalov determinant; the numerator $10 = 20/2$ arises from
the norm $\langle\Lambda|\Lambda\rangle = c(5c+22)/10$ and the
structure constant $c_{334}^2$. All inputs are rational functions
of $c$. The Arnold residue used here is the rational boundary residue
of the four-point affine screen; no non-unit iterated integral occurs.
Thus the motivic lift is the Tate-scalar element
$10/[c(5c+22)]\in\MZV^{\mot}_0\otimes\mathbb Q(c)$, and the period map
acts as the identity on it.
\end{proof}

\begin{remark}[\label{rem:large-c-normalization}Large-$c$ normalization]
The large-$c$ asymptotic $S_4^{\mot}(\Vir_c) \sim 2/c^2$, not
$2/(5c^2)$; this is a rational identity and receives no motivic
correction. The factor $5$ in $5c+22$ cancels against the numerator
$10=2\cdot5$ in the leading term.
\end{remark}

\subsection{$S_5^{\mot}(\Vir_c)$}

\begin{proposition}[\label{prop:s5-vir-mot}Motivic lift of
$S_5(\Vir_c)$]
\ClaimStatusProvedHere
The motivic lift of the weight-$5$ Virasoro shadow coefficient is
\[
 S_5^{\mot}(\Vir_c) \;=\; \frac{-48}{c^2 (5c + 22)}
 \;\in\; \MZV^{\mot}_{0} \otimes \mathbb{Q}(c)
 \;\subset\; \Fil^W_{\le 5}\MZV^{\mot} \otimes \mathbb{Q}(c).
\]
The external shadow degree is $5$; the Brown motivic weight is $0$.
\end{proposition}

\begin{proof}
The numerical value $S_5(\Vir_c) = -48/[c^2(5c+22)]$ is computed by
Wick contraction and agrees with the master-equation recursion:
\[
 S_5 \;=\; -\frac{1}{5c} \sum_{\substack{3 \leq j \leq k \\ j + k = 7}}
 \varepsilon(j, k)\, j k\, S_j S_k
 \;=\; -\frac{1}{5c}\bigl(2 \cdot 3 \cdot 4 \cdot S_3 S_4\bigr)
 \;=\; -\frac{24}{5c} \cdot 2 \cdot \frac{10}{c(5c+22)}
 \;=\; \frac{-48}{c^2(5c+22)}
\]
using $S_3 = 2$, $S_4 = 10/[c(5c+22)]$, and
$\varepsilon(3,4) = 1$.
The displayed recurrence is a finite rational product-sum on the
rational inputs $S_3$ and $S_4$. It contains no non-unit Arnold
iterated integral. Hence the motivic lift lies in the weight-$0$
summand $\MZV^{\mot}_0\otimes\mathbb Q(c)$.
\end{proof}

\subsection{$W_3$ at shadow degree $6$: period boundary}

\begin{proposition}[\label{prop:s6-w3-mot}Motivic lift of
$S_6(W_3)$: necessary period source]
\ClaimStatusConditional
For the principal $\cW_3$ algebra, the canonical local data determine
rational $T$-line and $W$-line collision coefficients:
\[
 S_4^T=\frac{10}{c(5c+22)},\qquad
 S_4^W=\frac{2560}{c(5c+22)^3}.
\]
A non-rational motivic contribution to the shadow-degree-$6$ coefficient
can occur only from a non-zero Arnold/KZ period pairing in the $WW$
channel that survives ordered-to-symmetric averaging. In that case it
has the filtered form
\[
 S_6^{\mot}(W_3)
 \;=\; S_6^{\rat}(W_3)
 \;+\;\sum_{1\le b\le 6}\gamma_{6,b}(c)\,m_b,
 \qquad
 m_b\in\MZV^{\mot}_b,\quad \gamma_{6,b}(c)\in\mathbb Q(c).
\]
The local sources used here do not prove that the coefficients of
$\zet^{\mot}(3)^2$ or $\zet^{\mot}(6)$ are non-zero.
\end{proposition}

\begin{proof}
The landscape census gives the principal $\cW_3$ $T$-line coefficient
and the separate $W$-line coefficient displayed above. It also records
the $WW$ collision channel with pole orders $\{6,4,3,2,1\}$ in the
OPE and $\{5,3,2,1\}$ after the $d\log$ shift. Thus the spin-$3$
sector is a genuine possible source of Arnold iterated integrals beyond
the Virasoro scalar lane.

Appendix~\ref{app:arnold-relations} supplies the obstruction: ordered
Stasheff boundary data do not by themselves imply all pairwise
Arnold/Kohno flatness relations, and the passage to symmetric data
requires a named compatibility calculation. Therefore the existence of
a spin-$3$ channel gives only the necessary source of a motivic period.
To prove a non-zero $\zet^{\mot}(3)^2$ or $\zet^{\mot}(6)$ coefficient
one must compute the actual Arnold/KZ pairing and check that it survives
the averaging map. That computation is not present in the local sources
used for this chapter.
\end{proof}

\begin{remark}[\label{rem:w3-first-zeta3}Virasoro scalar lane and
higher-spin period lane]
The Virasoro calculations in this chapter prove rationality only for
the scalar residue recurrences they display. The $\cW_3$ spin-$3$
channel is the first local place where a genuine Arnold/KZ period can
enter, because it supplies collision data not present on the Virasoro
$T$-line. Non-zero motivic content is a separate period-pairing
calculation, not a consequence of the existence of the spin-$3$
generator alone.
\end{remark}

%%%%%%%%%%%%%%%%%%%%%%%%%%%%%%%%%%%%%%%%%%%%%%%%%%%%%%%%%%%%%%%%%%%%%%%%%%%%%
\section{Virasoro rationality: weighted and manuscript towers}
\label{sec:virasoro-motivic-rationality-all-r}
%%%%%%%%%%%%%%%%%%%%%%%%%%%%%%%%%%%%%%%%%%%%%%%%%%%%%%%%%%%%%%%%%%%%%%%%%%%%%

The canonical sources separate two Virasoro sequences. The weighted
Riccati-metric coefficients $\widehat S^{\rm wt}_r$ satisfy the
one-line transport recurrence. The manuscript coefficients $S_r$ agree
with the weighted sequence through $r=5$ and split at $r=6$:
\[
 \widehat S^{\rm wt}_6(\Vir_c)
 =\frac{80(45c+193)}{3c^3(5c+22)^2},
 \qquad
 S_6(\Vir_c)=\frac{4(240c+1031)}{c^3(5c+22)^2}.
\]
The rationality theorem below is unconditional for the weighted tower
and records the proof boundary for the unweighted manuscript tower.

\begin{theorem}[\label{thm:virasoro-motivic-rationality-all-r}Virasoro
motivic rationality: weighted tower and manuscript boundary]
\ClaimStatusProvedHere
For every $r\ge2$, the weighted Riccati-metric Virasoro coefficient is
rational in the central charge:
\[
 \widehat S_{r}^{\rm wt,\mot}(\Vir_c)
 =
 \widehat S_{r}^{\rm wt}(\Vir_c)
 \in
 \MZV^{\mot}_{0}\otimes\mathbb Q(c)
 \subset
 \Fil^W_{\le r}\MZV^{\mot}\otimes\mathbb Q(c).
\]
Its motivic coaction is trivial:
\[
 \Delta^{\mot}\bigl(\widehat S_r^{\rm wt,\mot}(\Vir_c)\bigr)
 =
 1\otimes \widehat S_r^{\rm wt,\mot}(\Vir_c).
\]
The unweighted manuscript coefficients are rational in the range where
the null-state constraints have been computed, in particular through
the displayed $S_6$ above. An all-$r$ statement for the unweighted
sequence requires an all-$r$ rational null-state recurrence; it is not a
consequence of the weighted Riccati recurrence alone.
\end{theorem}

\begin{proof}
The weighted recurrence is
\[
 \widehat S^{\rm wt}_r(\Vir_c)
 =
 -\frac{1}{r c}
 \sum_{\substack{j+k=r+2\\3\le j\le k<r}}
 f(j,k)\,jk\,
 \widehat S^{\rm wt}_j(\Vir_c)\widehat S^{\rm wt}_k(\Vir_c),
 \qquad r\ge5,
\]
with $f(j,k)=1$ for $j<k$ and $f(j,j)=1/2$. This is the recurrence of
Theorem~\ref{thm:virasoro-shadow-recurrence}; the factor $r c$ is
$2r\kappa$ with $\kappa(\Vir_c)=c/2$. The initial values are
\[
 \widehat S^{\rm wt}_2=\frac c2,\qquad
 \widehat S^{\rm wt}_3=2,\qquad
 \widehat S^{\rm wt}_4=\frac{10}{c(5c+22)}.
\]
Strong induction gives $\widehat S^{\rm wt}_r(\Vir_c)\in\mathbb Q(c)$:
each summand is a rational scalar times a product of lower rational
functions, and division by $r c$ preserves $\mathbb Q(c)$.

The motivic lift of a rational function is its Tate-scalar element in
$\MZV^{\mot}_0\otimes\mathbb Q(c)$. Brown's coaction on the ground field
is $q\mapsto 1\otimes q$, so all positive-weight GRT components vanish
on these coefficients. The manuscript coefficient at $r=6$ is not
computed by this recurrence; the quadrichotomy chapter derives it from
the null-state constraint
$2S_2S_6+2S_3S_5+S_4^2=0$, giving the displayed rational value and
showing explicitly where the weighted and unweighted towers diverge.
\end{proof}

\begin{remark}[\label{rem:e1-vs-e2-rationality}The $E_1$-chiral residue
recurrence versus the $E_2$ Drinfeld associator]
Kontsevich's universal formality $\mathcal{U} \colon T_{\mathrm{poly}}
\to D_{\mathrm{poly}}$ involves motivic periods beyond
$\mathbb{Q}$: the depth-$2$ coefficient of Kontsevich's wheel graph
carries $\zet^{\mot}(3)$, and higher-depth coefficients populate
$\MZV^{\mot}$ (Rossi--Willwacher 2014, Alekseev--Torossian 2012). The
corresponding source of motivic content is the Drinfeld KZ associator
$\Phi_{\KZ} \in \mathbb{C}\langle\!\langle x, y\rangle\!\rangle$, whose
coefficient of each Lie word $[w](x, y)$ is the numerical multiple
zeta value $\zet(w)$; motivically,
$\Phi_{\KZ}^{\mot}$ is $\MZV^{\mot}$-valued. This generates the full
$\grt$-coaction on the $E_2$ Gerstenhaber operad.

Theorem~\ref{thm:virasoro-motivic-rationality-all-r} may appear
paradoxical next to this $E_2$-level richness: the Virasoro shadow
tower is $E_2$-visible (it is a coefficient of the universal
$\MC$-element of an $E_1$-chiral algebra whose derived center receives
an $E_2$-structure via Deligne--Tamarkin), so one might expect
$\MZV^{\mot}$-content to propagate from $\Phi_{\KZ}^{\mot}$ to the
Virasoro coefficients. The weighted Virasoro theorem does not make that
claim. It says that the one-dimensional $T$-line residue recurrence
uses only rational OPE constants and rational Arnold boundary residues.
The $E_2$ Drinfeld--Kontsevich periods belong to the associator and
derived-centre lane. They can affect a shadow coefficient only after a
specific Arnold/KZ period pairing survives the ordered-to-symmetric
averaging map. Proposition~\ref{prop:s6-w3-mot} isolates the first
higher-spin place where such a pairing may occur; it does not assert
non-vanishing without that computation.
\end{remark}

\begin{remark}[\label{rem:all-r-extends-r-leq-11}Relation to
Corollary~\ref{cor:virasoro-motivic-purity-r-leq-11}]
Corollary~\ref{cor:virasoro-motivic-purity-r-leq-11} belongs to the
weighted higher-coefficient chapter and is proved by explicit closed
forms through $r=11$. Theorem~\ref{thm:virasoro-motivic-rationality-all-r}
gives the all-$r$ induction for that weighted recurrence. The unweighted
manuscript sequence is governed by the separate null-state constraints
of the quadrichotomy chapter.
\end{remark}

%%%%%%%%%%%%%%%%%%%%%%%%%%%%%%%%%%%%%%%%%%%%%%%%%%%%%%%%%%%%%%%%%%%%%%%%%%%%%
\section{Kummer congruences from motivic Bernoulli-Kummer}
\label{sec:kummer}
%%%%%%%%%%%%%%%%%%%%%%%%%%%%%%%%%%%%%%%%%%%%%%%%%%%%%%%%%%%%%%%%%%%%%%%%%%%%%

The genus expansion of the characteristic-datum part
(\ref{part:characteristic-datum}) produces closed-form $Z_g(n)$ for
$g = 0,\ldots,9$. In the $\mathrm{SL}_2$ normalisation of
Chapter~\ref{ch:z-g-kummer-bernoulli}, the leading coefficient is
\[
 [n^{3g-3}]Z_g(k)
 =
 (-1)^g\frac{2^{g-1}B_{2g-2}}{(2g-2)!}.
\]
Thus the Bernoulli-leading irregular primes $691$ and $3617$ occur at
$g=7$ and $g=9$. The statement below is $p$-adic: it concerns the
Bernoulli numerator in the genus slot, not the Riccati characteristic
primes of the Virasoro shadow coefficients.

\begin{theorem}[\label{thm:kummer-from-motivic}Kummer congruences
from motivic Bernoulli-Kummer]
\ClaimStatusConditional
Let $p$ be an irregular prime with Bernoulli-irregularity index
$(2g-2, p)$ (i.e.\ $p \mid \mathrm{numer}(B_{2g-2})$). Assume the
motivic shadow-Feynman comparison identifies the motivic genus-$g$
slot with the $\mathrm{SL}_2$ Verlinde leading coefficient up to a
$p$-adic unit. If $p>2g-2$, then the genus slot has the same
$p$-adic numerator divisibility as
\[
 (-1)^g\frac{2^{g-1}B_{2g-2}}{(2g-2)!}.
\]
Consequently the Bernoulli-leading part of $Z_7$ contains the numerator
prime $691$, and the Bernoulli-leading part of $Z_9$ contains the
numerator prime $3617$.
\end{theorem}

\begin{proof}
The leading coefficient formula is
Theorem~\ref{thm:z-g-leading-coefficient-bernoulli}. Its scalar
multiplier $(-1)^g2^{g-1}/(2g-2)!$ is a $p$-adic unit whenever
$p>2g-2$. This holds for the first two relevant cases:
$691>12$ and $3617>16$.

Kummer's congruence
\[
 \frac{B_{2g-2}}{2g-2} \pmod{p} \text{ depends only on }
 (2g-2) \bmod (p-1)
\]
is a statement in the $p$-adic zeta/Tate lane; the comparison
hypothesis transports that $p$-adic class to the genus slot. The
$p$-adic-unit multiplier cannot remove a Bernoulli numerator prime.
For $g = 7$ ($2g-2 = 12$), $B_{12} = -691/2730$: the numerator
$691$ occurs in the leading coefficient. For $g = 9$ ($2g-2 = 16$),
$B_{16} = -3617/510$: the numerator $3617$ occurs in the leading
coefficient.
\end{proof}

\begin{remark}[\label{rem:kummer-verification}Kummer verification]
Direct computation of the $Z_g$ closed forms tests only the stated
Bernoulli-leading numerator divisibility. It does not imply that the
same primes occur in the Virasoro Riccati coefficients; the higher
coefficient chapter proves the opposite scope for the Bernoulli-leading
pair, with the isolated Laurent-tier $691$ occurrence at $\Gamma_8$.
\end{remark}

\begin{remark}[\label{rem:characteristic-primes-are-riccati-arithmetic}Characteristic
primes of the shadow tower are Riccati-recurrence integers]
\ClaimStatusProvedHere
The primes $\{61,\, 193,\, 2111,\, 16657,\, 14777,\, 3988097,\,
4969967,\, \ldots\}$ that appear in $S_r(\Vir_c)$ numerators
through $r = 11$ might be suspected of being Kac-module
characteristic primes: roots of the Shapovalov--Kac determinant
$\det G_N(c, h)$ at vacuum $h = 0$ for some level $N$. Direct
computation disproves this conjecture.

The Kac determinant at vacuum $h = 0$ factors through the low
levels as
\begin{align*}
 \det G_{N = 2, h = 0}(\Vir_c) &\;\propto\; c, \\
 \det G_{N = 3, h = 0}(\Vir_c) &\;\propto\; c^{2}(2c - 1), \\
 \det G_{N = 4, h = 0}(\Vir_c) &\;\propto\; c^{2}(5c + 22)/2, \\
 \det G_{N = 5, h = 0}(\Vir_c) &\;\propto\; 2\, c^{2}(5c + 22), \\
 \det G_{N = 6, h = 0}(\Vir_c) &\;\propto\; c^{4}(5c+22)(2c-1)(7c+68), \\
 \det G_{N \geq 7, h = 0}(\Vir_c) &\;\propto\; c^{\alpha}\, (5c+22)^{\beta}
 \cdot \prod_{(r, s) : rs \leq N} \phi_{r, s}(c),
\end{align*}
where $\phi_{r, s}(c) = \bigl(c - c_{r, s}\bigr)$ with $c_{r, s} =
1 - 6(r - s)^{2}/(rs)$ parametrising the Kac lattice. The
characteristic primes of Kac-module factors arise from the
$c_{r, s}$ rationals: $c_{1, 2} = -22/5, c_{1, 3} = -7,
c_{1, 4} = -46/7, c_{2, 3} = 1/2, c_{1, 5} = -22/5$, etc. These
rational factorisations yield the Kac-module characteristic-prime
set $\{2, 3, 5, 7, 11, 17, 22, 23, 46, 68, \ldots\}$; none of
these rationals generate $\{61, 193, 16657, 14777\}$.

Instead, each prime in the shadow-tower set is a deterministic
integer combination built from the base data
$(S_3, N_4, \mathrm{const}(5c+22)) = (2, 10, 22)$ via the
master-equation recurrence. Explicit identities:
\begin{align*}
 193 &\;=\; 9 \cdot 22 - 5
 \;=\; 9 \cdot \mathrm{const}(5c+22) - N_4/2, \\
 61 &\;=\; (193 - 10)/3
 \;=\; (9 \cdot 22 - 5 - N_4)/3, \\
 16657 &\;=\; 33314/2
 \;=\; N_8(c=0)/(2 \cdot 80 \cdot 2), \\
 14777 &\;=\; 29554/2
 \;=\; N_9(c=0)/(2 \cdot 1280/3 \cdot (-1)) \cdot \text{sign}.
\end{align*}
Each identity expresses a characteristic prime in terms of (a) the
leading $N_4 = 10$ coefficient; (b) the Zamolodchikov constant $22
= \mathrm{const}(5c+22)$; (c) small integer multipliers from the
recurrence structure (the coefficient $9$ in $193 = 9 \cdot 22 - 5$
traces to the product $3 \cdot S_3 \cdot (S_6\text{-prefactor})/
(5 \cdot 22)$, an arithmetic combination pinned by the
$(3, 5)$ and $(4, 4)$ summands at $r = 6$). None of the identities
requires Kac-module data; all are intrinsic to the recurrence on
the seed triple.

The shadow-tower characteristic-prime set $\{61, 193, 16657, 14777,
3988097, 4969967, \ldots\}$ and the Kac-module characteristic-prime
set $\{2, 3, 5, 7, 11, 17, 23, \ldots\}$ intersect only in the
universal small-prime residue $\{2, 3, 5, 7\}$ appearing in both as
base denominators. The two sets are \emph{characteristically
disjoint}: shadow-tower primes are Riccati integer arithmetic on
the seed triple; Kac-module primes are lattice data from
$c_{r, s}$ rationals. Thus the Kac-determinant explanation is false for
the displayed shadow primes.
\end{remark}

%%%%%%%%%%%%%%%%%%%%%%%%%%%%%%%%%%%%%%%%%%%%%%%%%%%%%%%%%%%%%%%%%%%%%%%%%%%%%
\section{Shadow L-function weight obstruction}
\label{sec:shadow-L}
%%%%%%%%%%%%%%%%%%%%%%%%%%%%%%%%%%%%%%%%%%%%%%%%%%%%%%%%%%%%%%%%%%%%%%%%%%%%%

The shadow $L$-function was introduced in the characteristic-datum part (\ref{part:characteristic-datum}) as the formal
Dirichlet series
\[
 L^{\mathrm{sh}}(\cA; s)
 \;=\; \sum_{r \geq 2} \frac{S_r(\cA)}{r^s}.
\]
The motivic lift supplies the weight filtration for this formal series. It
does not by itself prove analytic continuation or locate analytic
poles; those belong to the separate shadow-$L$ analysis.

\begin{corollary}[\label{cor:shadow-L-pole}Motivic weight obstruction
for the shadow $L$-function]
\ClaimStatusConditional
The formal motivic shadow $L$-series
\[
 L^{\mathrm{sh}, \mot}(\cA; s)
 \;=\; \sum_{r \geq 2} \frac{S_r^{\mot}(\cA)}{r^s}
\]
has no intrinsic weight-$1$ motivic residue: $\MZV^{\mot}_1=0$ and the
shadow tower starts at $r=2$. The coefficient $S_2=\kappa$ is a
Tate-scalar shadow invariant in
$\MZV^{\mot}_0\otimes\mathbb Q(c,k,\ldots)$ with external shadow
degree $2$, not a forced non-zero $\zet^{\mot}(2)$ period.
Consequently the bare rule
$F_g \leftrightarrow L^{\mathrm{sh}}(1-2g)$ is not supplied by the
motivic shadow tower. A motivic genus-slot comparison must add, as
separate data, an admissible analytic continuation, a coefficient
extraction or regularisation map, and the Tate twist matching the
Bernoulli-leading genus line.
\end{corollary}

\begin{proof}
Definition~\ref{def:mzv-mot} gives $\MZV^{\mot}_1=0$. Since the
series defining $L^{\mathrm{sh},\mot}$ begins at $r=2$, there is no
coefficient in shadow degree $1$. The coefficient in shadow degree $2$
is the conductor $\kappa$; for Virasoro it is $c/2$, a rational scalar.
By Theorem~\ref{thm:shadow-tower-motivic-lift}, its lift lies in the
weight-$0$ part of the filtered motivic coefficient ring unless a
separate Arnold period computation supplies a non-unit period.

The genus coefficient has leading Bernoulli factor
$B_{2g-2}/(2g-2)!$, hence Tate weight $2g-2$ in the $p$-adic Bernoulli
lane of Theorem~\ref{thm:kummer-from-motivic}. Evaluation of a
Dirichlet series at $s=1-2g$ is an analytic regularisation operation,
not coefficient extraction from the $r=2g-2$ shadow slot. Thus a
motivic identity at this point requires extra comparison data; the
formal motivic shadow tower alone only gives the filtered coefficient
sequence.
\end{proof}

%%%%%%%%%%%%%%%%%%%%%%%%%%%%%%%%%%%%%%%%%%%%%%%%%%%%%%%%%%%%%%%%%%%%%%%%%%%%%
\section{Shadow kappa and Dyson beta}
\label{sec:kappa-vs-beta}
%%%%%%%%%%%%%%%%%%%%%%%%%%%%%%%%%%%%%%%%%%%%%%%%%%%%%%%%%%%%%%%%%%%%%%%%%%%%%

The conductor $\kappa$ is the shadow coefficient $S_2$. Dyson's
$\beta$ is not a chiral-algebra invariant by itself; it is the
symmetry index of a chosen random-matrix ensemble. At $c=13$,
$\kappa(\Vir_{13})=13/2$, while the classical Dyson indices lie in
$\{1,2,4\}$. The equality problem is a type error: the two numbers
come from different structures.

\begin{theorem}[\label{thm:kappa-vs-beta-split}Motivic kappa,
modular beta]
\ClaimStatusProvedHere
Let $\cA$ be chirally Koszul in the standard landscape.
\begin{enumerate}[(i)]
\item The Koszul depth
$\kappa^{\mot}(\cA)=S_2^{\mot}(\cA)$ is a Tate-filtered shadow
invariant:
\[
\kappa^{\mot}(\cA)\in
\MZV^{\mot}_{0}\otimes\mathbb Q(c,k,\ldots)
\subset
\Fil^W_{\le2}\MZV^{\mot}\otimes\mathbb Q(c,k,\ldots),
\]
unless a separate degree-$2$ Arnold period is present. For the standard
scalar conductor formulas no such period occurs.
\item Dyson's $\beta$ is defined only after choosing an ensemble
$\mathcal E$ whose antiunitary symmetry type is specified. Then
$\beta(\mathcal E)\in\{1,2,4\}$ records the orthogonal, unitary, or
symplectic class. The modular character vector of $\cA$ may be used as
input for such an ensemble, but the bar complex and the character
vector do not canonically determine $\beta$.
\item There is no natural comparison map sending the shadow conductor
to Dyson's symmetry index. Numerical coincidences at special central
charges are not identities of invariants.
\end{enumerate}
\end{theorem}

\begin{proof}
\emph{(i)} The canonical conductor formulas are rational:
\[
\kappa(\Vir_c)=\frac c2,\qquad
\kappa(V_k(\fg))=\frac{\dim(\fg)(k+h^\vee)}{2h^\vee},\qquad
\kappa(\mathcal W_N)=c(H_N-1).
\]
Their motivic lifts are the same rational functions in the
weight-$0$ summand of the filtered motivic coefficient ring. Brown's
coaction on the ground field is $q\mapsto1\otimes q$.

\emph{(ii)} Dyson's $\beta$ is not defined by the bar-complex
construction; it is defined by a \emph{separate} datum: the symmetry
class of a Hermitian random-matrix ensemble. A modular $S$-matrix or a
Verlinde character vector can supply spectral input, but it does not
choose the real, complex, or quaternionic symmetry class. That choice
is external to the chiral algebra unless an ensemble construction is
part of the data.

\emph{(iii)} $\kappa$ is extracted from the ordered bar conductor.
$\beta$ is extracted from the symmetry type of a modular-character
ensemble. Regulators may compare selected motivic and automorphic
periods, but they do not identify a rational shadow conductor with a
Dyson symmetry class. At $c=13$ the values $13/2$ and
$\{1,2,4\}$ already lie in different numerical sets.
\end{proof}

\begin{remark}[\label{rem:kappa-beta-numerical-coincidence}Numerical
near-coincidence]
At the self-dual Virasoro point $c=13$, $\kappa(\Vir_{13})=6.5$.
Dyson's $\beta$ is an integer symmetry index. The fixed point of
$c\mapsto26-c$ is a shadow-conductor statement; the symmetry type of a
modular $S$-matrix is a separate modular statement.
\end{remark}

%%%%%%%%%%%%%%%%%%%%%%%%%%%%%%%%%%%%%%%%%%%%%%%%%%%%%%%%%%%%%%%%%%%%%%%%%%%%%
\section{Epilogue}
\label{sec:epilogue}
%%%%%%%%%%%%%%%%%%%%%%%%%%%%%%%%%%%%%%%%%%%%%%%%%%%%%%%%%%%%%%%%%%%%%%%%%%%%%

The shadow tower is arithmetic in two distinct senses. Its rational
Riccati coefficients are integer recursions in the OPE constants; its
non-rational coefficients, when present, are Brown--Goncharov periods
of Arnold/KZ iterated integrals. The motivic coaction of
Theorem~\ref{thm:grt-motivic-coaction} acts on the latter period
subspace and is trivial on the rational Virasoro scalar lane.

The Bernoulli-leading Kummer primes belong to the genus expansion
through Theorem~\ref{thm:kummer-from-motivic}. The Riccati primes of
the Virasoro shadow tower belong to a different arithmetic mechanism.
The shadow $L$-series and the Dyson $\beta$ comparison are likewise
weight and source tests: neither supplies an identity between
incompatible invariants.

The open calculation is explicit: compute the Arnold/KZ pairings for
higher-spin $\mathcal W_N$ channels and check their survival under
ordered-to-symmetric averaging. Only then does a motivic weight
signature distinguish two families beyond their rational shadow
projection.

%%%%%%%%%%%%%%%%%%%%%%%%%%%%%%%%%%%%%%%%%%%%%%%%%%%%%%%%%%%%%%%%%%%%%%%%%%%%%
\section{Nekrasov $\Omega$-background and external degree}
\label{sec:nekrasov-omega-motivic}
%%%%%%%%%%%%%%%%%%%%%%%%%%%%%%%%%%%%%%%%%%%%%%%%%%%%%%%%%%%%%%%%%%%%%%%%%%%%%

\providecommand{\ep}{\epsilon}

The $\Omega$-background supplies an external equivariant degree; it is
not, by itself, Brown motivic weight. For a $4d$ $\mathcal N=2$ theory
on $\mathbb R^4_{\ep_1,\ep_2}$, the Nekrasov partition function
$Z_{\mathrm{inst}}(\ep_1,\ep_2)$
(Nekrasov, \texttt{hep-th/0206161}; Nekrasov--Okounkov 2006) defines
the $\Omega$-deformed free energy
\[
 \mathcal F(\ep_1,\ep_2)
 \;=\;
 \ep_1\ep_2\,\log Z_{\mathrm{inst}}(\ep_1,\ep_2)
\]
with a Taylor expansion in $\ep_1,\ep_2$ whose coefficients are
equivariant integrals, rational functions, or periods according to the
theory and chamber \cite{NakajimaYoshioka2003,ErnstromNaganoNakajima2015}.
Only after an Arnold/KZ period comparison is constructed may such a
coefficient be lifted to
$\Fil^W_{\le n}\MZV^{\mot}$. The self-dual phase
$\ep_1=-\ep_2=\ep$ collapses the two-variable grading to the single
$\hbar$-grading via $\hbar^2=-\ep^2=\ep_1\ep_2$. Under this collapse
$\ep^n$ has external $\Omega$-degree $n$; a rational coefficient
carrying that monomial still has Brown weight~$0$. The comparison below
is therefore a conjectural match between shadow degree and external
$\Omega$-degree, with Brown weight supplied only by explicit
Arnold/KZ periods.

\begin{conjecture}[\label{conj:omega-motivic-weight-match}
$\Omega$-degree controls filtered shadow periods]
\ClaimStatusConjectured
Let $\cA$ be a $C_2$-cofinite non-logarithmic chiral algebra that
is attached by the Beem--Rastelli construction to a $4d$
$\mathcal{N}=2$ superconformal theory $\mathcal T$ on the
$\Omega$-background
$\mathbb{R}^4_{\ep_1,\ep_2}$. Then the shadow-tower coefficient
$S_r^{\mot}(\cA)\in\Fil^W_{\le r}\MZV^{\mot}$ coincides, after choosing
the same Arnold/KZ period coordinates and up to a rational
family-dependent prefactor, with the motivic lift of the external
$\Omega$-degree-$r$ coefficient of the Nekrasov--Okounkov refined free
energy at the self-dual phase:
\[
 S_r^{\mot}(\cA)
 \;\sim_{\mathbb{Q}(\mathrm{family})}\;
 \bigl[\,\ep^r\,\bigr]\,
 \mathcal F_{\mathrm{NO}}\bigl(\mathcal T;\,\ep,-\ep\bigr).
\]
The bracket $[\ep^r]$ denotes external equivariant degree, not
projection to Brown weight~$r$; if the coefficient is rational, its
motivic lift lies in $\MZV^{\mot}_0$.
\end{conjecture}

\begin{remark}[Nekrasov--Okounkov normalisation for
Conjecture~\ref{conj:omega-motivic-weight-match}]
\label{rem:omega-motivic-weight-match-strategy}
Nekrasov--Okounkov (2006) express the $\Omega$-deformed free energy via
equivariant integration over framed-instanton moduli; the resulting
formula evaluates to a partition sum with hook-length weights. Under
the AGT normalisation, when $\mathcal T$ is a class-$\mathcal S$
theory, the same instanton series is identified with a
Liouville/Toda conformal block. Mixed-Tate
Arnold/KZ coordinates give motivic lifts for the genus-zero collision
periods appearing in such blocks (Andr\'e~2009; Brown 2012), but the
statement that the $\ep$-degree-$r$ collision coefficient equals the
shadow coefficient is precisely Conjecture~\ref{conj:omega-motivic-weight-match}.
The self-dual phase $\ep_1=-\ep_2$ restricts to the single
$\hbar$-expansion used in the convolution dGLA. Family-dependent
rational normalisations, including the Beem--Rastelli dictionary
$c_{2d}=-12c_{4d}$, are absorbed in
$\sim_{\mathbb Q(\mathrm{family})}$. The conjecture is open beyond the
family-by-family checks for $\Vir_c$ (Liouville,
Nekrasov--Shatashvili 2009) and affine Kac--Moody
(Nekrasov--Pestun 2012).
\end{remark}

\begin{remark}[K3 instance: $\hbar^2 = -1/8$ vs.\ self-dual
$\ep^2 = 1/8$]
\label{rem:omega-k3-self-dual}
For the K3 chiral bialgebra $\mathbf{H}_{\Delta_5}$ of Vol~III, the
parameter $\hbar^2 = -1/8$ (Vol~III K3 Hall--Drinfeld pentagon-coboundary
calculation)
corresponds under $\hbar^2 = \ep_1\ep_2$ at the self-dual phase to
$\ep_1 = -\ep_2 = \pm 1/(2\sqrt{2})$ for the $\Omega$-deformed
class-$\mathcal S$ $A_1$ theory on $\Sigma_{0,24}$. The theorem
available at this point is the normalized AGT block comparison. The
Bethe algebra enters through the Nekrasov--Shatashvili degeneration,
which is adjacent to the fixed self-dual slice but not identical to it.
Identifying its Bethe wall crossing with Humbert walls, or the
Maulik--Okounkov stable-envelope $R$-matrix with a
Siegel-elliptic dynamical $R_{\mathrm{Sieg,dyn}}$, is a conditional
target comparison requiring the finite Schur--Igusa recognition data of
Remark~\ref{rem:feyn-nekrasov-schur-limit}
\cite{NekrasovShatashvili2009,MaulikOkounkov12}. This is the first
Vol~I/Vol~III test case for the Nekrasov external-degree comparison;
its generalisation to other CY-$2$ base inputs of the Vol~III
factorisation
$\Phi_d^{(\Sigma_{d-1},C)} = \SpCh_{\Sigma_{d-1}, C} \circ \PhiFA_d$ (the canonical
factorisation algebra $\PhiFA_d$ on the CY$_d$ base, followed by
specialisation) is open.
\end{remark}

\begin{remark}[Pentagon coboundary at $\hbar^3$: $25/3$ is not a
central charge]
\label{rem:omega-pentagon-coboundary}
The pentagon-equation coboundary at order $\hbar^3$ of the K3
Siegel--Borcherds associator decomposes as
\[
 \phi^{(3)}
 \;=\;
 \zeta(3)\,c_{\mathrm{symm}}
 \;+\;
 \tfrac{25}{3}\,c_{\mathrm{timelike}}
 \;+\;
 (\Phi_{10}/\eta^{24})\,c_{\Phi_{10}}
\]
(Vol~III K3 Hall--Drinfeld pentagon-coboundary calculation). The
coefficient $25/3$ equals
$\mathrm{rk}(\mathrm{II}_{25,1}) - 1$ divided by the triple-leg
antisymmetrisation denominator~$3$, i.e.\ the Fake Monster Cartan
rank minus the timelike Cartan direction. This is \emph{not} a
Virasoro central charge: there is no $c = 25/3$ minimal model
controlling this coefficient. The motivic weight-$3$
component $\zeta^{\mot}(3)\,c_{\mathrm{symm}}$ sits in $\MZV^{\mot}_3$,
while the $25/3$ coefficient is a weight-$0$ rational. Conflating
motivic weight-$3$ with a central-charge interpretation at the
pentagon coboundary is a type error;
primary-lit frame: Drinfeld 1990 quasi-Hopf, Etingof--Kazhdan
Parts~I--V.
\end{remark}

\section{Igusa--Dedekind factors}
\label{sec:mst-supplement}

\begin{remark}[Andrianov factorisation $L(s,\Phi_{10}) = L(s,\Delta\otimes E_6)\,\zeta(s-9)\,\zeta(s-8)$]
\label{rem:mst-1}
The automorphic factorisation used by the motivic shadow-tower
comparison is the Andrianov identity
\[
 L(s,\Phi_{10})
 \;=\;
 L(s,\Delta\otimes E_6)\,\zeta(s-9)\,\zeta(s-8),
\]
with $\Phi_{10}$ the Igusa cusp form of weight $10$ and
$\Delta\otimes E_6$ the convolution of the discriminant modular form
$\Delta\in S_{12}$ with the Eisenstein series $E_6$. The shifts $s-9$
and $s-8$ are pinned by the K3 Mukai signature: $9 = c_-/2 - 1 = 20/2 - 1$
and $8 = K^{\kappa_{\mathrm{ch}}} = 2c_+ = 8$ (Mukai doubling) \cite{Andrianov1979,GritsenkoNikulin1998}.
Cross-ref Vol.~III Chapter~\ref{V3-ch:k3-chiral-bialgebra-platonic}.
\end{remark}

\begin{remark}[Motivic fingerprint of the Igusa--Dedekind twist]
\label{rem:mst-2}
The period fingerprint of the
$\Phi_{10}^{\mathrm{un}}/\eta^{24}$ Borcherds twist is recorded by
the shifted pair of standard factors
$L(s,\Phi_{10}^{\mathrm{un}})$ and $L(s,\eta^{24})$
with $\eta^{24}$ the weight-$12$ cusp form on $\mathrm{SL}_2(\mathbb Z)$.
At the K3 base point, the $\eta^{24}$ factor encodes the $24$ Kodaira
$I_1$ fibres on $E^{\mathrm{nod,sm}}_{24}$ via the Dedekind denominator
on the $24$ nodal base, and the quotient lives on $\bar{\mathcal A}_2$
via Kuga--Satake. The motivic shadow tower thus organises as
$m_k(\Delta_5)$ coefficients reading the third-order pentagon
coboundary $\phi^{(3)}$ \cite{Borcherds1992,KugaSatake1967}. Cross-ref
Chapter~\ref{V3-ch:k3-chiral-bialgebra-platonic}.
\end{remark}

\section{Humbert residues and external $\Omega$-degree: the $A_1$
class-$\mathcal S$ Nekrasov realisation}
\label{sec:mst-nekrasov-A1-sigma-0-24}

Conjecture~\ref{conj:omega-motivic-weight-match} identifies the
external shadow degree of $S_r^{\mot}(\cA)$ with the $\ep^r$-degree of
the $\Omega$-deformed free energy; Brown weight is supplied only by the
period part of the coefficient.
At the K3 slice
$\hbar^2 = -1/8$ corresponds to the self-dual $\Omega$-point of the
class-$\mathcal S$ $A_1$ theory on $\Sigma_{0,24}$. The
three local comparisons are the normalized AGT Liouville block, the
finite Schur--Igusa target for Humbert residues, and the
Nekrasov--Shatashvili Bethe system as a cyclotomic comparison channel.
Only the first is a theorem at this site; the latter two require the
recognition hypotheses stated below.

\begin{proposition}[Humbert residues carry external $\Omega$-degree at
self-dual; \ClaimStatusConditional]
\label{prop:mst-nekrasov-humbert-motivic}
Let $\mathcal T = \mathcal T[A_1, \Sigma_{0,24}]$
be the class-$\mathcal S$ $A_1$ theory on the $24$-punctured sphere
with distinct marked points $z_1,\ldots,z_{24}$. The specialisation
$z_i=\zeta_{24}^{i-1}$ is a cyclotomic marking of the pointed sphere,
not a geometric Mathieu $M_{24}$ action on $\mathbb{CP}^1$. Let
$Z^{\mathrm{Nek}}_{\mathcal T}$ denote the $\Omega$-background
partition function evaluated at the self-dual slice
$\ep_1 = -\ep_2 = 1/(2\sqrt 2)$, so
$\hbar^2 = \ep_1\ep_2 = -1/8$.

Assume the finite Schur--Igusa recognition data of
Remark~\ref{rem:feyn-nekrasov-schur-limit}: the Beem--Rastelli
chiral algebra $\mathcal X(\mathcal T)$ is compared with
$\mathbf H_{\Delta_5}$ only after central-charge, level, root, parity,
and finite coefficient data identify the chosen K3-marked Schur
character with the reciprocal Igusa target; and the scalar Humbert-pole
comparison of Remark~\ref{rem:feyn-nekrasov-humbert-residues}
identifies
\[
 \mathrm{Res}_{H_n}\bigl[\log Z^{\mathrm{Nek}}_{\mathcal T}\bigr]
 \;=\; c_{\phi_{-2,1}}(n)\cdot [H_n].
\]
Then this conditional residue admits the motivic lift
\[
 \bigl(\mathrm{Res}_{H_n}\bigl[\log Z^{\mathrm{Nek}}_{\mathcal T}\bigr]\bigr)^{\mot}
 \;=\; c_n\cdot [H_n]^{\mot}
 \;\in\; \MZV^{\mot}_{0}
 \otimes \mathrm{CH}^1(\overline{\mathcal A}_2)^{\mot},
\]
where $c_n=c_{\phi_{-2,1}}(n)\in\mathbb Q$ is the Fourier coefficient
of $\phi_{-2,1}$ and $[H_n]^{\mot}$ is the motivic class of the Humbert
divisor in
$\mathrm{CH}^1(\overline{\mathcal A}_2)^{\mot}$. The index $n$ is the
Humbert and $\Omega$-degree under this target comparison; it is not the
Brown weight of the rational integer $c_n$. The period map recovers the
rational coefficient
$c_n = c_{\phi_{-2,1}}(n)$ of
Eichler--Zagier~$1985$ \S~$9$.

\emph{Proof.} The AGT input is exactly the normalized block theorem of
Proposition~\ref{prop:feyn-nekrasov-self-dual}:
$Z^{\mathrm{Nek}}_{\mathcal T}$ at the self-dual slice is the
Liouville conformal block on the same $24$-punctured sphere with
$c_{\mathrm{Liou}}=25$. This theorem supplies no Mathieu action on the
marked sphere and no identification of the Beem--Rastelli chiral
algebra with $\mathbf H_{\Delta_5}$.

Beem--Rastelli~$2013$ Thm.~$1$ identifies the Schur index with the
graded character of the associated $2d$ chiral algebra
$\mathcal X(\mathcal T)$. The passage from
$\mathrm{ch}_q(\mathcal X(\mathcal T))$ to the reciprocal Igusa
character is precisely the finite recognition hypothesis in the
statement, not a consequence of Beem--Rastelli. Under that hypothesis
the scalar Humbert-pole comparison reduces the local calculation to the
Siegel slice of $\log(\Phi_{10}^{\mathrm{un}}/\eta^{24})$. Its Fourier
expansion admits the Borcherds product form
$\Phi_{10}/\eta^{24} = \prod_{(m,n,r) > 0}(1 - p^m q^n y^r)^{
c_{\phi_{-2,1}}(4mn - r^2)}$ (Borcherds~$1998$ JAMS~$11$
Thm.~$10.1$). Taking $\log$ and collecting the $p^n$-coefficient
at the Humbert-$H_n$ slice yields the rational residue
$c_{\phi_{-2,1}}(n)\cdot [H_n]$ (Bruinier~$2002$ LNM~$1780$
Prop.~$5.1$). The motivic lift proceeds on each factor. On the
coefficient side, $\phi_{-2,1}$ is a weak Jacobi form arising
from the K3 elliptic genus $\chi(\mathrm K3; \tau, z) = 2\phi_{0,1}
- \tfrac{1}{2}(\wp + 2 E_2)\phi_{-2,1}$ (Eichler--Zagier~$1985$
Thm.~$9.3$); its Fourier coefficients lie in
$\mathbb Q\subset \MZV^{\num}$. Each rational Fourier coefficient has
the canonical Tate-scalar motivic lift
$c_{\phi_{-2,1}}(n)\cdot1\in\MZV^{\mot}_0$. The external label $n$ is
the Humbert label and, under
Conjecture~\ref{conj:omega-motivic-weight-match}, the $\ep^n$-degree;
it is not the Brown weight of the rational integer. On the divisor
side, the Humbert divisors $H_n$ are algebraic cycles on
$\overline{\mathcal A}_2$, so they admit motivic classes
$[H_n]^{\mot}\in \mathrm{CH}^1(\overline{\mathcal A}_2)^{\mot}$
(van der Geer~$1988$ Ch.~IX). The product is the motivic
residue in question; the period map recovers the numerical
identity by definition of $\per$ (Definition~\ref{def:period-map}).
\qed

\emph{Primary.} Alday--Gaiotto--Tachikawa~$2009$ Thm.~$1$;
Beem--Rastelli~$2013$ Thm.~$1$; Gritsenko--Nikulin~$1998$
Thm.~$2.1$; Borcherds~$1998$ JAMS~$11$ Thm.~$10.1$;
Bruinier~$2002$ LNM~$1780$ Prop.~$5.1$; Eichler--Zagier~$1985$
\S~$9$; van der Geer~$1988$ \emph{Hilbert Modular Surfaces}
Ch.~IX; Brown~$2012$ \emph{Ann.\ Math.}~$175$; see also
Remark~\ref{rem:feyn-nekrasov-schur-limit} and
Remark~\ref{rem:feyn-nekrasov-humbert-residues}.
\end{proposition}

\begin{corollary}[$c_3 = -8$ at $H_3$: explicit verification;
\ClaimStatusConditional]
\label{cor:mst-nekrasov-c3-H3}
Under the hypotheses of
Proposition~\ref{prop:mst-nekrasov-humbert-motivic}, at $n = 3$ the
motivic Humbert residue of
$\log Z^{\mathrm{Nek}}_{\mathcal T[A_1,\Sigma_{0,24}]}$ is
\[
 \bigl(\mathrm{Res}_{H_3}\bigl[\log Z^{\mathrm{Nek}}_{\mathcal T}\bigr]\bigr)^{\mot}
 \;=\; -8\cdot [H_3]^{\mot},
\]
whose period is the rational
$\mathrm{Res}_{H_3}\bigl[\log Z^{\mathrm{Nek}}\bigr] = -8\cdot [H_3]$
of the finite target comparison in Vol~I feynman\_diagrams.tex
Remark~\ref{rem:feyn-nekrasov-humbert-residues}.

\emph{Proof.} By Proposition~\ref{prop:mst-nekrasov-humbert-motivic}
the residue equals $c_{\phi_{-2,1}}(3)^{\mot}\cdot [H_3]^{\mot}$.
Direct $q, y$-expansion of the infinite-product form
$\phi_{-2,1} = -(y - 2 + y^{-1})\prod_{n\geq 1}(1 - q^n y)^2
(1 - q^n/y)^2/(1-q^n)^4$ through $q^8$ yields
$c_{\phi_{-2,1}}(3) = -8$, $c_{\phi_{-2,1}}(4) = +12$,
$c_{\phi_{-2,1}}(7) = -39$, $c_{\phi_{-2,1}}(8) = +56$ (Vol~I
bv\_brst.tex eq.~\eqref{eq:bvbrst-heegner-identity}; Eichler--Zagier
$1985$ \S~$9$ Table~$1$). At $n = 3$ this gives $c_3^{\mot} = -8$
as a $\mathbb Q$-rational Tate scalar, i.e.\ Brown weight~$0$ with
external Humbert/$\Omega$-degree~$3$; its period is
$-8 \in \mathbb Q\subset \mathbb R$. The admissibility condition
$n\equiv 0, 3\pmod 4$ (Eichler--Zagier~$1985$ Thm.~$2.2$) confirms
$n = 3$ is on the Fourier support. \qed
\end{corollary}

\begin{remark}[Cyclotomic Bethe-spectrum Galois shadow;
\ClaimStatusHeuristic]
\label{rem:mst-nekrasov-bethe-motivic}
The Nekrasov--Shatashvili limit of
$\mathcal T[A_1, \Sigma_{0,24}]$ identifies its free energy with the
Yang--Yang functional of the Gaudin model of type $A_1$ on
$\mathbb{CP}^1$ with $24$ marked points
(Feigin--Frenkel--Reshetikhin~$1994$ Comm.\ Math.\ Phys.~$166$;
Frenkel~$2007$ Ch.~$6$; Vol~I feynman\_diagrams.tex
Remark~\ref{rem:feyn-nekrasov-NS-bethe}). This is an NS degeneration,
not the fixed self-dual point $\ep_1=-\ep_2=1/(2\sqrt2)$.

The Bethe-ansatz spectrum $\{u_a\}$ is the zero set of a system of
$\bar{\mathbb Q}$-rational equations. Its algebraic coordinates are
acted on by $\mathrm{Gal}(\bar{\mathbb Q}/\mathbb Q)$, not by the
mixed-Tate motivic Galois group unless a separate Arnold/KZ period
torsor has been constructed. After the cyclotomic marking
$z_i=\zeta_{24}^{i-1}$, with equal external masses, the equations are
defined over $\mathbb Q(\zeta_{24})$ and acquire the evident
cyclotomic symmetry: rotation by $\zeta_{24}$ sends a solution
$\{u_a\}$ to $\{\zeta_{24}u_a\}$, while
$\mathrm{Gal}(\mathbb Q(\zeta_{24})/\mathbb Q)$ permutes the marked
points arithmetically. This is $\mu_{24}$ and cyclotomic Galois data,
not an $M_{24}$-equivariant system. A Mathieu-moonshine interpretation
would require an additional comparison map from this cyclotomic Gaudin
system to the Cheng--Duncan--Harvey K3 module.

Under the finite Igusa-recognition hypotheses of
Proposition~\ref{prop:mst-nekrasov-humbert-motivic}, the Humbert
residue target in the Bethe channel is the absolute-Galois module
generated by
$\{c_{\phi_{-2,1}}(n)\cdot [H_n]^{\mot}\}_{n\geq3,\,
n\equiv0,3(4)}$, with $c_{\phi_{-2,1}}(n)$ a Tate scalar. A genuine
$\mathrm{Gal}^{\mot}(\mathrm{MT}(\mathbb Z))$-comodule would require
non-rational KZ/Arnold periods in the Bethe integral. Its period, if
such a torsor is constructed, is the numerical Bethe spectrum of the
$24$-point Gaudin system in the NS degeneration. Primary:
Feigin--Frenkel--Reshetikhin~$1994$; Frenkel~$2007$ Ch.~$6$;
Nekrasov--Shatashvili~$2009$ arXiv:$0908.4052$; Brown~$2012$
\emph{Ann.\ Math.}~$175$ for the motivic Galois framework;
Cheng--Duncan--Harvey~$2014$ for the moonshine target.
\end{remark}

%%%%%%%%%%%%%%%%%%%%%%%%%%%%%%%%%%%%%%%%%%%%%%%%%%%%%%%%%%%%%%%%%%%%%%%%%%%%%
\section{The motivic bar--cobar bridge on $\overline{\mathcal A}_{2}$}
\label{sec:motivic-barcobar-bridge}
%%%%%%%%%%%%%%%%%%%%%%%%%%%%%%%%%%%%%%%%%%%%%%%%%%%%%%%%%%%%%%%%%%%%%%%%%%%%%

\providecommand{\BettiRlz}{\mathrm{Betti}}
\providecommand{\HeegObs}{\mathrm{Heeg}}
\providecommand{\Data}{\mathrm{Data}}
\providecommand{\Heeg}{\mathrm{Heeg}}
\providecommand{\MT}{\mathrm{MT}}

The chain-level compatibility $\infty$-groupoid
$\Data(\cA)\simeq B\mathrm{Aut}^h_{L_\infty}(\mathfrak g^{\infty,1}_{\cA})$
of Theorem~B is a torsor under $\widehat{\mathrm{grt}}_1$ modulo the
Gaitsgory--Nadler obstruction $\mathrm{ob}^{\mathrm{GN}}$
(\cite{GaitsgoryNadler09}; Vol~I
\texttt{e3\_identification\_chain\_level\_platonic.tex}
Chapter~\ref{chap:bar-cobar-adjunction}). The Siegel--Eisenstein pairing
$\langle [\chi_3], [e_3^{\mathrm{K3}\times E}]\rangle_{\Phi_3}
= 2\,\mathrm{Vol}(E)(2\pi i)^3
= \chi(\cO_{\mathrm{K3}})\cdot\Res_{s=1/2}[E_5^{(2)}(Z,s)]|_{K(1)}$
of Theorem~H expresses a numerical period. The bridge conjecture asserts
that both data are images under a single period realisation of a common
motivic object, and that this compatibility holds away from the Humbert
locus where Heegner--Chern torsion obstructs.

\begin{definition}[Motivic compatibility-data groupoid;
\ClaimStatusConjectured]
\label{def:motivic-data-groupoid}
Let $\mathscr D$ be a finite set of admissible Humbert discriminants
and set
\[
 U_{\mathscr D}
 :=
 \overline{\mathcal A}_{2,\mathbb Q}
 \setminus \bigcup_{D\in\mathscr D} H_D .
\]
This is the algebraic open on which the following object is defined.
The phrase ``away from all Humbert divisors'' means the inverse system
over all finite $\mathscr D$, not a single finite-type open complement
of the infinite Humbert union. Let
$\Data^{\mot}_{\mathscr D}(\cA)$ denote the conjectural pro-object in
$\mathcal{DM}(U_{\mathscr D};\mathbb Q)$ obtained from the motivic
homotopy automorphisms of a relative ordered bar complex
$B_{U_{\mathscr D}/\mathbb Q}^{\mathrm{ord},\mot}(\cA)$ over
$U_{\mathscr D}$:
\[
 \Data^{\mot}_{\mathscr D}(\cA)
 \;:=\;
 \mathbf{R}\mathrm{Aut}_{\mathcal{DM}(U_{\mathscr D})}^{h,L_\infty}\!
 \bigl(\mathfrak g^{\infty,1,\mot}_{\cA,U_{\mathscr D}}\bigr).
\]
The motivic Galois action on the mixed-Tate Arnold/KZ torsor maps,
through the Ihara--Drinfeld action, to its image in
$\widehat{\mathrm{grt}}_1$. No quotient
$\mathrm{Gal}^{\mot}(\MT(\mathbb Z))\twoheadrightarrow\grt$ is asserted.
The period realisation is required to give a comparison morphism
\[
 \per_{\mathscr D}\colon \Data^{\mot}_{\mathscr D}(\cA)
 \longrightarrow
 \Data^{\BettiRlz}(\cA)|_{U_{\mathscr D}}
 \;=\; \Data(\cA)_{\CC}|_{U_{\mathscr D}}
\]
to the Betti compatibility groupoid of Theorem~B tensored to
$\mathbb C$.
\end{definition}

\begin{definition}[Heegner--Chern obstruction module;
\ClaimStatusConditional]
\label{def:heeg-chern-obstruction}
For each admissible discriminant $D\in\{1, 4, 5, 8, 9, 12, 13, \ldots\}$
(discriminants $D\equiv 0, 1\bmod 4$) let
$\mathrm{cl}(H_D)\in H^2_{\mathrm{mot}}(\overline{\mathcal A}_{2};
\mathbb Q(1))$ be the motivic cycle class of the Humbert divisor
$H_D$ (Borcherds 1998; Kudla--Millson 1990). The cycle classes are
unconditional. The obstruction module additionally requires a chosen
action of the motivic image in $\widehat{\mathrm{grt}}_1$ on their
rational span. For a finite set $\mathscr D$ as in
Definition~\ref{def:motivic-data-groupoid}, set
\[
 \Heeg_{\mathscr D} \;:=\;
 \bigoplus_{D\in\mathscr D} \mathbb Q\cdot \mathrm{cl}(H_D)
 \;\subset\;
 H^2_{\mathrm{mot}}(\overline{\mathcal A}_{2}; \mathbb Q(1))
\]
as a module over the motivic image. The candidate obstruction to
extending $\per_{\mathscr D}$ across the finite boundary
$\bigcup_{D\in\mathscr D}H_D$ is the class
\[
 \mathrm{ob}^{\mot}_{\Data,\mathscr D}(\cA)
 \;\in\;
 H^2\bigl(\mathfrak g_{\mathrm{mot}\to\mathrm{grt}};
 \Heeg_{\mathscr D}\bigr),
\]
where $\mathfrak g_{\mathrm{mot}\to\mathrm{grt}}$ denotes the image of
the motivic Lie algebra under the Ihara--Drinfeld/Willwacher comparison.
The all-Humbert obstruction is the compatible family of these finite
classes as $\mathscr D$ varies.
\end{definition}

\begin{conjecture}[\label{thm:motivic-barcobar-bridge}Motivic bar--cobar
bridge]
\ClaimStatusConjectured
Fix $\cA = \mathbf H_{\Delta_5}$, the K3-fibre BKM crown of Vol~III.
The period realisation of
Definition~\ref{def:motivic-data-groupoid} restricts, for every finite
admissible $\mathscr D$, to an equivalence
\[
 \per_{\mathscr D}\colon
 \Data^{\mot}_{\mathscr D}(\cA)
 \xrightarrow{\;\sim\;}
 \Data^{\BettiRlz}(\cA)|_{U_{\mathscr D}}
\]
over $U_{\mathscr D}$, compatibly under enlarging $\mathscr D$. The
Siegel--Eisenstein pairing of Theorem~H lifts to a motivic pairing
\[
 \bigl\langle [\chi_3]^{\mot}, [e_3^{\mathrm{K3}\times E}]^{\mot}
 \bigr\rangle_{\Phi_3}^{\mot}
 \;=\;
 2\,\mathrm{Vol}(E)^{\mot}\cdot(2\pi i)^{\mot,3}
 \;\in\;
 \Fil^W_{\le3}\MZV^{\mot}
 \otimes_{\mathbb Q} H^0_{\mathrm{mot}}(E; \mathbb Q(0))
\]
whose weight-$3$ projection has period recovering the
numerical identity
$2\,\mathrm{Vol}(E)(2\pi i)^3$ of the preface
eq.~\eqref{eq:bvbrst-heegner-identity} only after the Heegner
obstruction classes vanish:
\[
 \mathrm{ob}^{\mot}_{\Data,\mathscr D}(\cA) \;=\; 0
 \text{ in }
 H^2(\mathfrak g_{\mathrm{mot}\to\mathrm{grt}};\,\Heeg_{\mathscr D})
 \qquad\text{for every finite }\mathscr D .
\]
\end{conjecture}

\begin{remark}[Conditional construction of
Conjecture~\ref{thm:motivic-barcobar-bridge} under the
Deligne--Goncharov/Voevodsky/Ayoub framework]
\label{rem:motivic-bridge-evidence}
\emph{Motivic Arnold piece over $U_{\mathscr D}$.}
Deligne--Goncharov 2005 Theorem~5.25 lifts the genus-zero KZ/Arnold
iterated integrals on $\overline{M}_{0,n}$ to mixed-Tate motivic
periods. For $B^{\mathrm{ord}}(\cA)$ over
$U_{\mathscr D}\subset\overline{\mathcal A}_{2,\mathbb Q}$ this
supplies the coefficient periods on affine collision screens. A
relative object
$B^{\mathrm{ord},\mot}_{U_{\mathscr D}/\mathbb Q}(\cA)
\in\mathcal{DM}(U_{\mathscr D})$ carrying the OPE residue maps is the
extra datum demanded by
Definition~\ref{def:motivic-data-groupoid}; it is not a consequence of
the Deligne--Goncharov theorem alone. Its de~Rham and Betti
realisations are required to be the chain-level ordered bar complex and
the corresponding singular-chain complex.

\emph{Bar--cobar inversion, Koszul duality, and the bulk.}
Assuming this relative motivic bar complex and the usual Koszul-locus
conservativity for its realisations, Theorem~A lifts to a motivic
adjunction $\Omega^{\mot}\dashv B^{\mot}$ on the generated Tate
subcategory. Its period realisation is the bar--cobar inversion
$\Omega_XB_X(\cA)\to\cA$ in $D^{\mathrm{co}}_{\mathrm{ch}}(X)$. This
object is not the Koszul dual algebra. The Koszul dual is
$\cA^!=(\cA^i)^\vee$ after Verdier duality on the bar coalgebra
$B(\cA)$; the chiral derived centre
$Z^{\mathrm{der}}_{\mathrm{ch}}(\cA)$ is the Hochschild closed-sector object of
Theorem~H. The motivic bridge compares period realisations of these
typed objects; it does not identify them.

\emph{Galois compatibility of $\Data^{\mot}$.}
The motivic Galois action on the Arnold/KZ iterated-integral torsor
maps, via Ihara--Drinfeld and Willwacher's graph-complex realisation, to
the $\widehat{\mathrm{grt}}_1$-torsor of compatibility $2$-morphisms in
Theorem~B (Willwacher 2015 Thm.~1.1; Brown 2012 Thm.~1.2). The
comparison is through the motivic image in $\widehat{\mathrm{grt}}_1$,
not through an equality of motivic Galois and Grothendieck--Teichmuller
Lie algebras. Equivariance of each $\per_{\mathscr D}$ is therefore
one of the conditions in
Conjecture~\ref{thm:motivic-barcobar-bridge}.

\emph{Motivic lift of the Eisenstein pairing.}
The Chern class part of
$[e_3^{\mathrm{K3}\times E}]$ has a motivic lift by the cycle-class map.
A motivic lift of $[\chi_3]$ is additional Hochschild-regulator data:
it must realise the curve-level chiral Hochschild class of Theorem~H and
pair with the Chern class under Beilinson's regulator formalism
(Beilinson 1984 \emph{Higher regulators}). Under this hypothesis, the
de~Rham period of the pairing is the regularised Siegel--Eisenstein
residue of Kudla--Rallis 1994.

\emph{Heegner obstruction at the Humbert boundary.}
For finite $\mathscr D$, the boundary of
$U_{\mathscr D}$ in $\overline{\mathcal A}_{2,\mathbb Q}$ is the
finite union of Humbert divisors $H_D$, $D\in\mathscr D$. Their motivic
cycle classes form $\Heeg_{\mathscr D}$. Extending
$\Data^{\mot}_{\mathscr D}$ across this boundary would require a
splitting of the extension
\[
 0\to
 \Data^{\mot}_{\mathscr D}
 \to
 \widetilde{\Data}^{\mot}_{\mathscr D}
 \to
 \bigoplus_{D\in\mathscr D}
 \mathrm{cl}(H_D)\otimes_{\mathbb Q}\Heeg_{\mathscr D}
 \to 0,
\]
and the obstruction is the class of
Definition~\ref{def:heeg-chern-obstruction}. The codimension-$3$
calculation at $(D_1,D_2,D_3)=(3,4,7)$ verifies the numerical
Siegel--Eisenstein period identity on that locus. Global extension
across the all-Humbert pro-system requires the compatible vanishing of
all finite obstruction classes
\(\mathrm{ob}^{\mot}_{\Data,\mathscr D}\).
\end{remark}

\begin{conjecture}[\label{cor:motivic-lift-chi3}Motivic form of the
$\chi_3$-pairing]
\ClaimStatusConjectured
If
\[
 \mathrm{ob}^{\mot}_{\Data,\mathscr D}(\mathbf H_{\Delta_5})=0
 \qquad\text{for every finite admissible }\mathscr D,
\]
then
\[
 \bigl\langle [\chi_3]^{\mot}, [e_3^{\mathrm{K3}\times E}]^{\mot}
 \bigr\rangle_{\Phi_3}^{\mot}
 \;=\; \chi(\cO_{\mathrm{K3}})^{\mot}\cdot
 \Res^{\mot}_{s=1/2}\!\bigl[E_5^{(2),\mot}(Z, s)\bigr]\big|_{K(1)}
\]
in the weight-$3$ Eisenstein period line inside
$\Fil^W_{\le3}\MZV^{\mot}$. Here $E_5^{(2),\mot}(Z,s)$ and
$\Res^{\mot}_{s=1/2}$ denote a chosen motivic lift of the
regularised Siegel--Eisenstein residue on the paramodular group $K(1)$;
existence and compatibility of this lift are part of the conjecture.
Period realisation recovers the numerical identity of the preface.
\end{conjecture}

\begin{remark}[Three compatibility tests for
Conjecture~\ref{thm:motivic-barcobar-bridge}]
\label{rem:motivic-bridge-three-paths}
The conjecture has three independent compatibility tests. Agreement of
the tests is evidence; it does not prove the motivic equivalence
$\per_{\mathscr D}$ for all finite $\mathscr D$.

\textup{(i)} \emph{Deligne--Goncharov motivic $\pi_1$.} Apply the
Deligne--Goncharov 2005 motivic iterated-integral construction to the
affine collision screens on a generic fibre of
$\overline{M}_{0,2g+2}$ over $U_{\mathscr D}$. The expected comparison
is that the
resulting motivic $\pi_1^{\mot}$-torsor realises the same associator
coordinates as the $\widehat{\mathrm{grt}}_1$-torsor of chain-level
$\Data(\cA)$, with Furusho 2011 providing associator rigidity. This
test is internal to mixed Tate motives over~$\mathbb Z$ and does not
use Siegel--Eisenstein.

\textup{(ii)} \emph{Ikeda lift of Siegel--Eisenstein.} Ikeda 2001
(\emph{Invent.\ Math.}\ 144) supplies the automorphic lift controlling
the Siegel--Eisenstein residue. A motivic refinement of that residue,
when constructed, would lie on the weight-$3$ Eisenstein extension line
whose period is the classical Zagier value. This test is automorphic
and disjoint from the operadic~(i).

\textup{(iii)} \emph{K3-chiral Chern character.} The Mukai vector
of the K3 fibre at an admissible Humbert point maps to its motivic
Chern character $\mathrm{ch}^{\mot}\in
K_0^{\mot}(\mathrm{K3})\otimes_{\mathbb Z}\mathbb Q$; its pairing
with the degree-$3$ BKM class
$e_3^{\mathrm{K3}\times E}$ computes
$\chi(\cO_{\mathrm{K3}}) = 2$ and $\mathrm{Vol}(E)$ from the
fibre integral
(Mukai 1987; Oberdieck 2019 reduced DT; Vol~III K3 Hall--Drinfeld
Schur-index calculation).
This path is geometric and disjoint from both (i) and (ii).

Agreement across the three tests at codim-$3$ is verified
numerically at $(D_1, D_2, D_3) = (3, 4, 7)$ with
value $2\,\mathrm{Vol}(E)(2\pi i)^3$. Brown's description of the
weight-$3$ line gives a distinguished generator
$\zet^{\mot}(3)$; after fixing this normalisation the calculation
selects the corresponding motivic lift. Surjectivity of $\per$ alone
would not give uniqueness.
\end{remark}

\begin{remark}[Motivic and chain-level compatibility
data; \ClaimStatusConditional]
\label{rem:motivic-vs-chainlevel-data}
The compatibility data $\Data(\cA)$ admits two genuinely distinct
enhancements, never to be conflated:
\emph{chain-level} $\Data(\cA)\simeq
B\mathrm{Aut}^h_{L_\infty}(\mathfrak g^{\infty,1}_{\cA})$, a
$\widehat{\mathrm{grt}}_1$-torsor valued in simplicial sets over
$\mathbb C$ (Theorem~B), versus \emph{motivic}
$\{\Data^{\mot}_{\mathscr D}(\cA)\}_{\mathscr D}$, conjectural torsors
for the motivic image in $\widehat{\mathrm{grt}}_1$ over the finite
Humbert-complement opens
$U_{\mathscr D}\subset\overline{\mathcal A}_{2,\mathbb Q}$
(Definition~\ref{def:motivic-data-groupoid}). Period realisations
$\per_{\mathscr D}$ are conjectured to be equivalences compatibly in
$\mathscr D$; extension across Humbert divisors is obstructed by the
finite classes $\mathrm{ob}^{\mot}_{\Data,\mathscr D}$ in the
image-Lie-algebra cohomology of
Definition~\ref{def:heeg-chern-obstruction}
(Conjecture~\ref{thm:motivic-barcobar-bridge}). Conflation produces
the error of treating the Siegel--Eisenstein residue
$\Res_{s=1/2}[E_5^{(2)}]|_{K(1)}$ as a universal invariant of
$\Data(\cA)$ rather than a period of
$\Data^{\mot}_{\mathscr D}(\cA)$ at a specific Heegner datum; the numerical
coincidence $2\,\mathrm{Vol}(E)(2\pi i)^3$ is a \emph{restricted}
shadow, not the whole motivic structure.
\end{remark}

%%%%%%%%%%%%%%%%%%%%%%%%%%%%%%%%%%%%%%%%%%%%%%%%%%%%%%%%%%%%%%%%%%%%%%%%%%%%%

\begin{remark}[Nekrasov motivic scope]
\label{rem:mst-nekrasov-scope}
Four distinctions are load-bearing. (i) The coefficient
$c_n$ of $\phi_{-2,1}$ is weight-$0$ on the $\MZV^{\num}$ side
(a rational integer), while the monomial carrying it has external
$\ep^n$-degree in Conjecture~\ref{conj:omega-motivic-weight-match}.
(ii) $\hbar^2 = -1/8$ is \emph{not} a central charge:
it is the product $\ep_1\ep_2$ at the self-dual slice; conflating
it with a Virasoro central charge is a type error. (iii) The Beem--Rastelli
identification is a \emph{character} trace, not an algebra
isomorphism; see Vol~III toric\_cy3\_coha.tex
Remark~\ref{V3-rem:K3xE-coha-chiral-guard}. (iv) The cyclotomic
marking $z_i=\zeta_{24}^{i-1}$ is not a geometric $M_{24}$ action on
the $24$-punctured sphere; Moonshine enters only through a separate K3
module comparison.
\end{remark}

%%%%%%%%%%%%%%%%%%%%%%%%%%%%%%%%%%%%%%%%%%%%%%%%%%%%%%%%%%%%%%%%%%%%%%%%%%%%%
\section{Single-valued scope for chiral-Hochschild periods}
\label{sec:sv-scope-restriction}
%%%%%%%%%%%%%%%%%%%%%%%%%%%%%%%%%%%%%%%%%%%%%%%%%%%%%%%%%%%%%%%%%%%%%%%%%%%%%

\providecommand{\zetasv}{\zeta^{\mathrm{sv}}}
\providecommand{\MZVsv}{\MZV^{\mathrm{sv}}}
\providecommand{\MZVmot}{\MZV^{\mot}}
\providecommand{\proj}{\mathrm{pr}^{\mathrm{sv}}}

Deligne--Goncharov attach motivic periods to the KZ/Arnold
iterated-integral torsor. Brown's single-valued period map sends
motivic multiple zeta values to the real-analytic quotient
$\MZVsv$. The chiral-Hochschild pairing of Theorem~H uses real Arnold
forms on the real Fulton--MacPherson compactification. That reality
condition is the source of the single-valued restriction.

\begin{theorem}[\label{thm:sv-scope-restriction-chiralhoch}Single-valued
scope of the chiral-Hochschild period identity;
\ClaimStatusConditional]
Assume the motivic Hochschild pairing of
Conjecture~\ref{cor:motivic-lift-chi3}
\[
 \bigl\langle [\chi_r]^{\mot},
 [\omega_r^{\mathrm{univ}}]^{\mot}\bigr\rangle_{\mathrm{cyc}}^{\mot}
 \;\in\; \Fil^W_{\le r}\MZVmot
\]
is represented by real Arnold forms
$d\log|z_i-z_j|^2$ on
$\overline{\mathrm{Conf}_n(X)}(\RR)$. Then its real period factors
through Brown's single-valued projection
$\proj\colon \MZVmot\to\MZVsv$. Equivalently, for each finite
$\mathscr D$ the relevant period-realisation map
$\per_{\mathscr D}$ of Definition~\ref{def:motivic-data-groupoid}
factors as
\[
 \per_{\mathscr D}\colon
 \Data^{\mot}_{\mathscr D}(\cA)\;\xrightarrow{\;\proj_*\;}\;
 \Data^{\mathrm{sv}}(\cA)\;\xrightarrow{\;\mathrm{real}\;}\;
 \Data^{\BettiRlz}(\cA)|_{U_{\mathscr D},\RR},
\]
and the image of the Hochschild pairing in
$\Data^{\mathrm{sv}}$ is a $\mathrm{Gal}(\CC/\RR)$-invariant
motivic real period. The Siegel--Eisenstein numerical identity
$2\,\mathrm{Vol}(E)(2\pi i)^3 = \chi(\cO_{\mathrm{K3}})
\Res_{s=1/2}[E_5^{(2)}(Z,s)]|_{K(1)}$ of Theorem~H is a period
in the single-valued realisation of the weight-$3$ period line, not a
generic element carrying the full motivic Galois action.
\end{theorem}

\begin{proof}
Arnold harmonic representatives $\omega_r^{\mathrm{univ}}$ are real
forms on $\overline{\mathrm{Conf}_n(X)}(\RR)$: the Arnold algebra
is generated by
$\eta_{ij} = d\log|z_i - z_j|^2 = (dz_{ij}/z_{ij} + \overline{dz_{ij}/z_{ij}})$,
which is the single-valued real form whose holomorphic part is
the KZ one-form. Integration of a motivic iterated integral
against a single-valued harmonic form lands in $\MZVsv$ by
Brown 2013 \S 2.1 Thm.~2.4 (single-valued period realisation). The
Hochschild cocycle $\chi_r$ is, under the hypothesis of the theorem, the
cyclic-Arnold pairing
$\langle\chi_r,\omega_r^{\mathrm{univ}}\rangle_{\mathrm{cyc}}$ by
the shadow-ChirHoch bridge, which is a real-period pairing. The factorisation
$\per_{\mathscr D} = \mathrm{real}\circ\proj_*$ follows from the
Brown 2013 single-valued factorisation diagram of the motivic
comparison isomorphism on each finite open. The Siegel--Eisenstein
residue at $s=1/2$ of
$E_5^{(2)}$ is a regularised real period
(Kudla--Rallis 1994 regularised Siegel--Weil), so any motivic lift of
this identity lands in the single-valued realisation. In weight $3$ the
single-valued odd-zeta line is generated by $\zetasv(3)$, with Brown's
normalisation fixing the scalar.
\end{proof}

\begin{remark}[Scope separations for
Theorem~\ref{thm:sv-scope-restriction-chiralhoch}]
\label{rem:sv-scope-five-obstructions}
Five separations are essential. First, $S_2=c/2$ is a rational shadow
coefficient, not a depth-one MZV; the vanishing of weight-one
single-valued MZVs does not constrain it. Second, Heegner--Chern
classes control extension of the motivic domain across Humbert
divisors, while the single-valued projection controls the codomain of
real periods. Third, the Zamolodchikov norm $c(5c+22)/10$ is a rational
normalisation, not an iterated-integral period. Fourth, the
Gangl--Kaneko--Zagier double-zeta relations are compatible with
single-valued projection after taking the real invariant quotient.
Fifth, Kontsevich graph weights live in the full motivic period algebra;
only the real Arnold-Hochschild subclass is forced into the
single-valued image.
\end{remark}

\begin{remark}[Motivic, single-valued, and chain-level sites;
single-valued scope; \ClaimStatusConditional]
\label{rem:motivic-sv-crosslevel}
The compatibility data $\Data(\cA)$ thus inhabits \emph{three}
distinct sites, never to be conflated.
\textup{(a)} Chain-level
$\Data(\cA)\simeq B\mathrm{Aut}^h_{L_\infty}(\mathfrak g^{\infty,1}_{\cA})$
over $\mathbb C$, a $\widehat{\mathrm{grt}}_1$-torsor of
simplicial sets (Theorem~B).
\textup{(b)} Motivic
$\{\Data^{\mot}_{\mathscr D}(\cA)\}_{\mathscr D}$, a conjectural
compatible system of torsors for the motivic image in
$\widehat{\mathrm{grt}}_1$ over the finite Humbert-complement opens
$U_{\mathscr D}\subset\overline{\mathcal A}_{2,\mathbb Q}$
(Definition~\ref{def:motivic-data-groupoid}).
\textup{(c)} Single-valued $\Data^{\mathrm{sv}}(\cA)$, the image of
$\Data^{\mot}_{\mathscr D}(\cA)$ under Brown's single-valued period
projection, carrying
$\mathrm{Gal}(\CC/\RR)$-invariant real motivic periods
(Theorem~\ref{thm:sv-scope-restriction-chiralhoch}).
The real-period refinement of the chiral-Hochschild pairing lives in
(c); Theorem~H itself is a curve-level Hochschild statement in (a).
Under the hypotheses of Theorem~\ref{thm:sv-scope-restriction-chiralhoch},
period realisation factors as
$(b)\xrightarrow{\proj_*}(c)\xrightarrow{\mathrm{real}}(a)_\RR$.
Vol~III's factorisation
$\Phi_d^{(\Sigma_{d-1},C)} = \SpCh_{\Sigma_{d-1}, C} \circ \PhiFA_d$
combines the canonical $E_d$-holomorphic factorisation algebra
$\PhiFA_d$ with specialisation
$\SpCh_{\Sigma_{d-1}, C}$ to the reference curve $C$ by
factorisation homology over a $(d{-}1)$-cycle. The BKM denominator
$\chi_{10}$ is a real-analytic Borcherds product
(Borcherds 1992 Thm.~10.1), so its real periods are compared through
the single-valued image of $\MZVmot$, not through arbitrary motivic
coordinates. Theorem~H's concentration bound
$\mathrm{ChirHoch}^\bullet\subset\{0,1,2\}$ is a curve-level
Hochschild theorem; the single-valued restriction is a period-realisation
refinement, not a replacement for the chain-level proof.
\end{remark}
